% 前置内容：在这里填写论文信息、中英文摘要和关键词。
% 由main.tex加载，不单独编译；排版设置仍在geshi.tex。
% ==================== 01 论文信息 ==================== 中文标题正标题不超过25字，正副标题合计不超过35字。
\newcommand{\thesisTitle}{在此填写中文论文题目}
\newcommand{\thesisTitleEn}{YOUR THESIS TITLE IN ENGLISH}
\newcommand{\thesisType}{学术学位论文} % 专硕改为：专业学位论文
\newcommand{\authorName}{填写姓名}
\newcommand{\authorNameEn}{Your Name}
\newcommand{\supervisor}{填写导师及职称}
\newcommand{\supervisorEn}{Prof. Your Supervisor}
\newcommand{\degreeName}{工学硕士}
\newcommand{\degreeNameEn}{Master of Engineering}
\newcommand{\specialty}{建筑学}
\newcommand{\specialtyEn}{Architecture}
\newcommand{\submissionDate}{2026年12月}
\newcommand{\defenseDate}{2026年12月}
\newcommand{\defenseDateEn}{December, 2026}
\newcommand{\classification}{待填写}
\newcommand{\udc}{待填写}
\newcommand{\schoolCode}{待填写}


% ==================== 02 封面与内封 ====================
\makecover\makeinnertitles\frontmatter
% ==================== 03 中英文摘要 ====================
\begin{cnabstract}
在此填写中文摘要。依次概括研究背景、研究问题、研究内容与方法、主要结果和结论。摘要应独立完整，与英文摘要对应。
\keywords{关键词一；关键词二；关键词三}
\end{cnabstract}
\begin{enabstract}
Write the English abstract here. Its content should correspond to the Chinese abstract.
\enkeywords{keyword one, keyword two, keyword three}
\end{enabstract}

\clearpage\tableofcontents

